\documentclass[../../../include/open-logic-section]{subfiles}

\begin{document}

\olfileid[tr]{sth}{ordinals}{basic}
\olsection{Sıra Sayılarının Temel Özellikleri}

İlk birkaç sıra sayısının doğal sayılar olduğunu gözlemledik. Sıra sayıları
teorisi geliştirmenin başlıca nedeni, doğal sayılarda geçerli tümevarım
ilkesini genişletmektir. Buna bir dizi temel sonuçla adım adım ulaşacağız.

\begin{lem}\ollabel{ordmemberord}
Bir sıra sayısında bulunan her !!{element}, bir sıra sayısıdır.
\end{lem}

\begin{proof}
$\alpha$ bir sıra sayısı ve $b\in\alpha$ olsun. $\alpha$ geçişli olduğundan
$b\subseteq\alpha$ olur. Dolayısıyla $\in$, $\alpha$'yı iyi sıraladığı gibi
$b$'yi de iyi sıralar.

$b$'nin geçişli olduğunu görmek için $x\in c\in b$ varsayın.
$b\subseteq\alpha$ olduğundan $c\in\alpha$'dır. Yine $\alpha$ geçişli
olduğundan $c\subseteq\alpha$ ve dolayısıyla $x\in\alpha$ olur. Böylece
$x,c,b\in\alpha$'dır. Fakat $\in$, $\alpha$'yı iyi sıralar; bu yüzden
\olref[sth][ordinals][wo]{wo:strictorder} gereği $\in$, $\alpha$ üzerinde
geçişli bir bağıntıdır. $x\in c\in b$ olduğundan $x\in b$ elde edilir.
Genelleyerek $c\subseteq b$ sonucuna ulaşırız.
\end{proof}

\begin{cor}\ollabel{ordissetofsmallerord}
Her sıra sayısı $\alpha$ için
$\alpha=\Setabs{\beta\in\alpha}{\beta\text{ bir sıra sayısıdır}}$.
\end{cor}

\begin{proof}
\olref{ordmemberord} sonucundan doğrudan çıkar.
\end{proof}

Sonraki iki ana sonucun, \olref{ordinductionschema} ile
\olref{ordtrichotomy}'nin, kabaca söylediği şey sıra sayılarının
kendilerinin üyelikle iyi sıralandığıdır:

\begin{thm}[Sonlu Ötesi Tümevarım]\ollabel{ordinductionschema}
Her $\phi(x)$ formülü için:
\[
	\text{eğer }\exists \alpha \phi(\alpha)\text{ ise }\exists \alpha(\phi(\alpha)
	\land  (\forall \beta \in \alpha) \lnot \phi(\beta))
\]
olur; gösterilen niceleyiciler örtük olarak sıra sayılarıyla sınırlıdır.
\end{thm}

\begin{proof}
%\olref{ordinductionschema1}. 
Bir sıra sayısı $\alpha$ için $\phi(\alpha)$ olduğunu varsayın. Eğer
$(\forall\beta\in\alpha)\lnot\phi(\beta)$ ise işimiz biter. Değilse,
$\alpha$ bir sıra sayısı olduğundan $\phi$'yi sağlayan bir $\in$-en küçük
!!{element} içerir; bu eleman da \olref{ordmemberord} gereği bir sıra
sayısıdır.
%	\olref{ordinductionschema2}. Suppose there is some ordinal
%	$\gamma$ such that $\lnot\phi(\gamma)$. Then by
%	\olref{ordinductionschema1}, there is an $\in$-minimal ordinal
%	$\alpha$ for which $\lnot\phi(\alpha)$. So $(\forall \beta <
%	\alpha) \phi(\beta)$, rendering the antecedent of the conditional
%	false.
\end{proof}
\noindent
\olref{ordinductionschema} sonucunu, \olref{ordinductionschema} içinde
$\lnot\phi(\alpha)$ alıp temel mantıksal dönüşümleri yaparak, eşdeğer biçimde
şu şema olarak ifade edebiliriz:
\[
\text{eğer }\forall \alpha((\forall \beta \in \alpha)\phi(\beta) \lif
\phi(\alpha))\text{ ise }\forall \alpha\phi(\alpha).
\]

\begin{thm}[Üçlü Karşılaştırma]\ollabel{ordtrichotomy}
Her $\alpha$ ve $\beta$ sıra sayısı için
$\alpha\in\beta\lor\alpha=\beta\lor\beta\in\alpha$.
\end{thm}

\begin{proof}
İspat çift tümevarımla, yani \olref{ordinductionschema} iki kez kullanılarak
yapılır. $x\in y\lor x=y\lor y\in x$ ise $x$'in $y$ ile
\emph{karşılaştırılabilir} olduğunu söyleyelim.

Tümevarım için $\alpha$ içindeki her sıra sayısının \emph{her} sıra
sayısıyla karşılaştırılabilir olduğunu varsayın. İkinci tümevarım için
$\alpha$'nın $\beta$ içindeki her sıra sayısıyla karşılaştırılabilir
olduğunu varsayın. $\alpha$'nın $\beta$ ile karşılaştırılabilir olduğunu
göstereceğiz. $\beta$ üzerinde tümevarımla $\alpha$'nın her sıra sayısıyla;
ardından $\alpha$ üzerinde tümevarımla \emph{her} sıra sayısının \emph{her}
sıra sayısıyla karşılaştırılabilir olduğu çıkar. Gereken budur.
$\alpha\notin\beta$ ve $\beta\notin\alpha$ varsayıp $\alpha=\beta$
olduğunu göstermek yeterlidir.

$\alpha\subseteq\beta$ olduğunu göstermek için $\gamma\in\alpha$
sabitleyin; \olref{ordmemberord} gereği bu bir sıra sayısıdır. İlk tümevarım
varsayımına göre $\gamma$, $\beta$ ile karşılaştırılabilir. Fakat
$\gamma=\beta$ ya da $\beta\in\gamma$ olsaydı (gerektiğinde $\alpha$'nın
geçişli olduğunu kullanarak) $\beta\in\alpha$ çıkardı; bu varsayımımıza
aykırıdır. Öyleyse $\gamma\in\beta$ olur. Genelleyerek
$\alpha\subseteq\beta$ elde ederiz.

İkinci tümevarım varsayımını kullanan bütünüyle benzer bir gerekçe
$\beta\subseteq\alpha$ olduğunu gösterir. Böylece $\alpha=\beta$ olur.
\end{proof}\noindent
Bu nedenle $\in$ bir sıralama bağıntısı gibi davrandığından bazen
$\alpha\in\beta$ yerine $\alpha<\beta$ yazacağız. Bunun alışkanlık dışında
derin bir nedeni yoktur; ayrıca $\alpha\in\beta\lor\alpha=\beta$ yerine
$\alpha\leq\beta$ yazmak daha kolaydır.\footnote{$\alpha
\mathrel{\underline{\in}}\beta$ yazabilirdik; fakat bu bütünüyle standart
dışı olurdu.}

Son sonuçlarımızın iki hızlı sonucu şöyledir; bunların ilki yeni gösterimimizi
kullanır:

\begin{cor}\ollabel{ordordered}
$\exists\alpha\phi(\alpha)$ ise
$\exists\alpha(\phi(\alpha)\land\forall\beta(\phi(\beta)\lif
\alpha\leq\beta))$ olur. Ayrıca her $\alpha,\beta,\gamma$ sıra sayısı için
hem $\alpha\notin\alpha$ hem de
$\alpha\in\beta\in\gamma\lif\alpha\in\gamma$ olur.
\end{cor}

\begin{proof}
Tıpkı \olref[wo]{wo:strictorder} gibidir.
\end{proof}

\begin{prob}
\olref[sth][ordinals][basic]{ordtrichotomy} ispatındaki ``bütünüyle benzer
gerekçeyi'' tamamlayın.
\end{prob}

\begin{cor}\ollabel{corordtransitiveord}
$A$, ancak ve ancak geçişli bir sıra sayıları kümesiyse bir sıra sayısıdır.
\end{cor}

\begin{proof}
\emph{Soldan sağa.} \olref{ordmemberord} gereği. \emph{Sağdan sola.}
$A$ geçişli bir sıra sayıları kümesiyse \olref{ordinductionschema} ile
\olref{ordtrichotomy} gereği $\in$, $A$'yı iyi sıralar.
\end{proof}

\olref{ordinductionschema} ile \olref{ordtrichotomy}'yi, sıra sayılarının
$\in$ ile iyi sıralandığını söyleyen sonuçlar olarak açıkladık. Ne var ki
aşağıdaki sonuç nedeniyle bu tür bir iddia karşısında \emph{çok dikkatli}
olmalıyız:

\begin{thm}[Burali--Forti Paradoksu]\ollabel{buraliforti}
Bütün sıra sayılarının bir kümesi yoktur.
\end{thm}

\begin{proof}
Olmayana ergiyle $O$'nun bütün sıra sayılarının kümesi olduğunu varsayın.
$\alpha\in\beta\in O$ ise \olref{ordmemberord} gereği $\alpha$ bir sıra
sayısıdır; dolayısıyla $\alpha\in O$ olur. Öyleyse $O$ geçişlidir ve
\olref{corordtransitiveord} gereği bir sıra sayısıdır. Böylece $O\in O$
olur; bu da \olref{ordordered} ile çelişir.
\end{proof}
\noindent
Bu sonuç adını \citeauthor{Burali-Forti1897}'den alır. Ancak bütün sıra
sayılarının bir kümesi bulunduğunu varsaymaktaki \emph{çelişkiyi} ilk kez
1899'da, Dedekind'e yazdığı bir mektupta Cantor açıkça görmüştür. van
Heijenoort şöyle açıklar:
\begin{quote}
  Burali--Forti, çelişkinin kendisini, olmayana ergi yoluyla, sıra
  sayılarının doğal sıralamasının yalnızca kısmi bir sıralama olduğu
  sonucunu kuruyor saydı.
  \citep[s.~105]{Heijenoort1967}
\end{quote}
Burali--Forti'nin hatasını bir yana bırakırsak yukarıdakileri şöyle
özetleyebiliriz: sıra sayıları üyelikle tek tek iyi sıralanmış ve üyelikle
topluca iyi sıralanmış kümelerdir (ama birlikte bir küme oluşturmazlar).

Son olarak, sıra sayıları hakkında \olref{ordinductionschema} ile
\olref{ordtrichotomy}'den çıkan birkaç temel özellik daha verelim.

\begin{prop}
Sıra sayılarından oluşan her kesin azalan dizi sonludur.
\end{prop}

\begin{proof}
Sıra sayılarından oluşan sonsuz bir kesin azalan
$\alpha_0>\alpha_1>\alpha_2>\ldots$ dizisinin $<$-minimal elemanı yoktur;
bu da \olref{ordinductionschema} ile çelişir.
\end{proof}

\begin{prop}\ollabel{ordinalsaresubsets}
Her $\alpha,\beta$ sıra sayısı için
$\alpha\subseteq\beta\lor\beta\subseteq\alpha$.
\end{prop}

\begin{proof}
$\alpha\in\beta$ ise $\beta$ geçişli olduğundan
$\alpha\subseteq\beta$ olur. Benzer biçimde $\beta\in\alpha$ ise
$\beta\subseteq\alpha$ olur. $\alpha=\beta$ ise hem
$\alpha\subseteq\beta$ hem $\beta\subseteq\alpha$ olur. Dolayısıyla
\olref{ordtrichotomy} gereği işimiz biter.
\end{proof}

\begin{prop}\ollabel{ordisoidentity}
Her $\alpha,\beta$ sıra sayısı için $\alpha=\beta$ ancak ve ancak
$\ordeq{\alpha}{\beta}$.
\end{prop}

\begin{proof}
Sıra sayıları iyi sıralamalardır; dolayısıyla sonuç Üçlü Karşılaştırma
(\olref[basic]{ordtrichotomy}) ile
\olref[iso]{wellordnotinitial} sonucundan doğrudan çıkar.
\end{proof}

\begin{prob}\ollabel{probunionordinalsordinal}
$X$'in her elemanı bir sıra sayısıysa $\bigcup X$'in bir sıra sayısı
olduğunu ispatlayın.
\end{prob}

\end{document}
